% ch04_u01 —— 原书 p.122–p.131（4.0 全部；4.1 自 4.1.1 起，至 p.131 末句中断处）
\bdchap{ch:pfaff}{Pfaffian 及其相关}
\SourcePage{122}

\bdsec{4.0. 引言}

\bdp{4.0.1} 考虑「归一化」的典范丛
\begin{equation}\tag{146}
\omega^{\sharp}_{\operatorname{Bun}_G}:=\omega_{\operatorname{Bun}_G}\otimes\omega_0^{\otimes -1},
\end{equation}
其中 $\omega_0$ 是 $\omega_{\operatorname{Bun}_G}$ 在 $\operatorname{Bun}_G$ 中对应于 $X$ 上平凡 $G$-丛那一点处的纤维。本节将给每个 ${}^{L}G$-oper $\mathfrak{F}$ 配置 $\operatorname{Bun}_G$ 上的可逆层 $\lambda_{\mathfrak{F}}$，即 0.2(d) 中提到的那一个。层 $\lambda_{\mathfrak{F}}$ 上会装一个同构 $\lambda_{\mathfrak{F}}^{\otimes 2n}\xrightarrow{\ \sim\ }(\omega^{\sharp}_{\operatorname{Bun}_G})^{\otimes n}$，其中 $n\ne 0$。该同构诱导出 0.2(d) 所要求的、$\lambda_{\mathfrak{F}}$ 上的扭 $D$-模结构。

由 3.4.3 的公式 (57)，$\operatorname{Op}_{{}^{L}G}(X)=\operatorname{Op}_{{}^{L}\mathfrak{g}}(X)\times Z\operatorname{tors}_{\theta}(X)$，其中 $Z$ 是 ${}^{L}G$ 的中心。实际上 $\lambda_{\mathfrak{F}}$ 只依赖 $\mathfrak{F}$ 在 $Z\operatorname{tors}_{\theta}(X)$ 中的像。于是目标就是构造一个典范函子
\begin{equation}\tag{147}
\lambda:Z\operatorname{tors}_{\theta}(X)\to\mu_\infty\operatorname{tors}_{\theta}(\operatorname{Bun}_G),
\end{equation}
其中 $\mu_\infty\operatorname{tors}_{\theta}(\operatorname{Bun}_G)$ 是如下群胚：对象为 $\operatorname{Bun}_G$ 上的线丛 $\mathcal{A}$，连同同构 $\mathcal{A}^{\otimes 2n}\xrightarrow{\ \sim\ }(\omega^{\sharp}_{\operatorname{Bun}_G})^{\otimes n}$（对某个 $n\ne 0$）。

\bdp{4.0.2} 若 $G$ 单连通，(147) 的构造相当简单。此时 $Z$ 平凡，只需构造 $\mu_\infty\operatorname{tors}_{\theta}(\operatorname{Bun}_G)$ 的一个对象。由于 $G$ 单连通，$\operatorname{Bun}_G$ 连通且单连通（把 $X$ 上的 $G$-丛解释为 $X$ 对应 $C^\infty$ 流形上的 $G$-丛，并带上一个 $\bar\partial$-联络）。于是问题化为说明 $\omega^{\sharp}_{\operatorname{Bun}_G}$ 有平方根（此时 $\mu_\infty\operatorname{tors}_{\theta}(\operatorname{Bun}_G)$ 有唯一对象，该对象在 $\operatorname{Bun}_G$ 中对应于平凡 $G$-丛那一点处的纤维被平凡化）。为解决此问题，我们使用 Pfaffian（普法夫式）的概念。

对带非退化对称型 $\mathcal{Q}\otimes\mathcal{Q}\to\omega_X$ 的向量丛 $\mathcal{Q}$，Laszlo 与 Sorger 在 [La-So] 中给 $\mathcal{Q}$ 配置其 Pfaffian $\operatorname{Pf}(\mathcal{Q})$，即 $\det R\Gamma(X,\mathcal{Q})$ 的典范平方根。4.2 中我们给出 Pfaffian 的另一个定义，它应当与 [La-So] 中的定义等价。

\SourcePage{123}

固定 $\mathcal{L}\in\omega^{1/2}(X)$（即 $\mathcal{L}$ 是 $\omega_X$ 的一个平方根）。则 $\operatorname{Bun}_G$ 上以
\begin{equation}\tag{148}
\operatorname{Pf}(\mathfrak{g}_{\mathfrak{F}}\otimes\mathcal{L})\otimes\operatorname{Pf}(\mathfrak{g}\otimes\mathcal{L})^{\otimes -1}
\end{equation}
为 $\mathfrak{F}\in\operatorname{Bun}_G$ 处纤维的线丛，是 $\omega^{\sharp}_{\operatorname{Bun}_G}$ 的一个平方根（细节见 4.3.1）。

于是要理解 $G$ 单连通的情形，只需通读 4.2 与 4.3.1。一般情形下 (147) 的构造更复杂。要点在于：(148) 所定义的 $\omega^{\sharp}_{\operatorname{Bun}_G}$ 的平方根依赖 $\mathcal{L}\in\omega^{1/2}(X)$。

\bdp{4.0.3} 下面概述对任意半单 $G$ 构造 (148) 的思路。

如 3.4.6 所述，$Z\operatorname{tors}_{\theta}(X)$ 是 Picard 范畴 $Z\operatorname{tors}(X)$ 上的一个 torsor，而 $\mu_\infty\operatorname{tors}_{\theta}(\operatorname{Bun}_G)$ 是 Picard 范畴
\begin{equation}\tag{149}
\mu_\infty\operatorname{tors}(\operatorname{Bun}_G):=\lim_{\substack{\longrightarrow\\ n}}\mu_n\operatorname{tors}(\operatorname{Bun}_G)
\end{equation}
上的一个 torsor。

% [Tutor Note: 原书 p.123 式 (149) 的 lim 下方印的是右向箭头（正极限的写法），但 μ∞-torsor 应为 μ_n-torsor 的逆极限；此处照录原书箭头，未改。]

我们要构造的函子 (147)，对某个 Picard 函子 $\ell:Z\operatorname{tors}(X)\to\mu_\infty\operatorname{tors}(\operatorname{Bun}_G)$ 而言，是 3.4.6 意义下的 $\ell$-仿射函子。$\ell$ 在 4.1 中定义。torsor $Z\operatorname{tors}_{\theta}(X)$ 由 $\omega^{1/2}(X)$ 经某个 Picard 函子 $\mu_2\operatorname{tors}(X)\to Z\operatorname{tors}(X)$ 诱导而来（见 3.4.6）。因此要构造 $\lambda$，只需构造一个 $\ell'$-仿射函子 $\lambda':\omega^{1/2}(X)\to\mu_\infty\operatorname{tors}_{\theta}(X)$，其中 $\ell'$ 是复合 $\mu_2\operatorname{tors}(X)\to Z\operatorname{tors}(X)\xrightarrow{\ \ell\ }\mu_\infty\operatorname{tors}(\operatorname{Bun}_G)$。我们以 $\mathcal{L}\mapsto\lambda'_{\mathcal{L}}$ 定义 $\lambda'$，其中 $\lambda'_{\mathcal{L}}$ 是 $\operatorname{Bun}_G$ 上的线丛，它在 $\mathfrak{F}\in\operatorname{Bun}_G$ 处的纤维等于 (148)。$\lambda'$ 是 $\ell'$-仿射这一事实，在 4.4 中由 4.3.10 推出；4.3.10 是关于 $\mathrm{SO}_n$-丛的一般论断\footnote{事实上，4.3.10 是 [BLaSo] 中命题 5.2 的精细化。}。实际上，在专门讨论 Pfaffian 的 4.2 与 4.3 两小节里，群 $G$ 完全不出现。

\bdp{4.0.4} 本节构造的 $\operatorname{Bun}_G$ 上每个线丛都带有如下附加结构：对每个 $x\in X$，$G(K_x)$ 的某个中心扩张作用在它到 2.3.1 中概形 $\operatorname{Bun}_{G,x}$ 的拉回上。这一结构在 4.3 中用到，第 5 章还会用到。

\SourcePage{124}

% 节标题进 \addcontentsline/书签：unicode-math 的希腊字母在 PDF 字符串里会触发
% 「Improper alphabetic constant」，故只对 $\mu_\infty$ 加 \texorpdfstring，正文排版不变。
\bdsec{4.1. $\operatorname{Bun}_G$ 上的 \texorpdfstring{$\mu_\infty$}{μ∞}-torsor}

\bdp{4.1.1} 设 $G$ 是连通仿射代数群，$\Pi$ 是有限交换群，$0\to\Pi(1)\to\tilde G\to G\to 0$ 是 $G$ 的一个扩张。我们要构造典范 Picard 函子 $\ell:\Pi^{\vee}\operatorname{tors}(X)\to\mu_\infty\operatorname{tors}(\operatorname{Bun}_G)$，其中 $\Pi^{\vee}:=\operatorname{Hom}(\Pi,\mu_\infty)$。

\emph{注.} 若 $G$ 半单、$\tilde G$ 是 $G$ 的泛覆叠，则 $\Pi=\pi_1(G)$，而 $\Pi^{\vee}$ 典范同构于 ${}^{L}G$ 的中心 $Z$（该同构由 $G$ 与 ${}^{L}G$ 的嘉当环面之间的对偶诱导）。故此时 $\ell$ 是 Picard 函子 $Z\operatorname{tors}(X)\to\mu_\infty\operatorname{tors}(\operatorname{Bun}_G)$，正如 4.0.3 所许诺的。

我们在 4.1.2–4.1.4 中构造 $\ell$。4.1.5 对上述构造给出「解释」，4.1.6 稍作改写。4.1.7–4.1.9 考虑 $G(K_x)$ 的某个中心扩张的作用。4.1.10–4.1.11 描述 $\operatorname{Bun}_G$ 的基本群胚，这使 $\operatorname{Bun}_G$ 上 torsor 的构造变得清楚。读者可以略过 4.1.5 与 4.1.10–4.1.11。

\bdp{4.1.2} 对 $\mathfrak{F}\in\operatorname{Bun}_G$，记 $\tilde{\mathfrak{F}}$ 为 $X$ 上 $\mathfrak{F}$ 的 $\tilde G$-提升构成的 $\Pi(1)$-gerbe。其类 $c(\mathfrak{F})$ 是 $cl(\mathfrak{F})$ 在边缘态射 $H^1(X,G)\to H^2(X,\Pi(1))=\Pi$ 下的像。对有限非空子集 $S\subset X$，gerbe $\tilde{\mathfrak{F}}_{X\setminus S}$ 平凡。因此 $|\tilde{\mathfrak{F}}(X\setminus S)|$（即 $\tilde{\mathfrak{F}}(X\setminus S)$ 的对象的同构类集合）是非空的 $H^1(X\setminus S,\Pi(1))$-torsor。记之为 $\phi_{S,\mathfrak{F}}$。$\mathfrak{F}$ 变动时，$\phi_{S,\mathfrak{F}}$ 成为 $\operatorname{Bun}_G$ 上某个 $H^1(X\setminus S,\Pi(1))$-torsor $\phi_S$ 的纤维。

\bdp{4.1.3} 对任意 $x\in X$，集合 $|\tilde{\mathfrak{F}}(\operatorname{Spec} O_x)|$ 只有一个元素。用它可平凡化 $\Pi$-torsor $|\tilde{\mathfrak{F}}(\operatorname{Spec} K_x)|$（注意 $\Pi=H^1(\operatorname{Spec} K_x,\Pi(1))$）。于是限制到 $\operatorname{Spec} K_s$（$s\in S$）给出 $\operatorname{Res}_s$-仿射映射 $\operatorname{Res}_{s,\mathfrak{F}}:\phi_{S,\mathfrak{F}}\to\Pi$，其中 $\operatorname{Res}_s:H^1(X\setminus S,\Pi(1))\to\Pi$ 是 $s$ 处的留数。对 $c\in\Pi$，令 $\Pi^S_c:=\{\pi_S=(\pi_s):\sum\pi_s=c\}\subset\Pi^S$。映射 $\operatorname{Res}_{S,\mathfrak{F}}:=(\operatorname{Res}_{s,\mathfrak{F}}):\phi_{S,\mathfrak{F}}\to\Pi^S$ 的像是 $\Pi^S_{c(\mathfrak{F})}$。

\SourcePage{125}

\bdp{4.1.4} 回忆 $\Pi^{\vee}$ 是对偶于 $\Pi$ 的群，故有非退化配对 $(\ ,\ ):\Pi\times\Pi^{\vee}\to\mu_\infty$。

设 $\mathcal{E}$ 是 $X$ 上的 $\Pi^{\vee}$-torsor。置 $\mathcal{E}_S:=\prod_{s\in S}\mathcal{E}_s$，即 $\mathcal{E}$ 在 $S$ 处的平凡化集合；它是 $(\Pi^{\vee})^S$-torsor。对任意 $e\in\mathcal{E}_S$ 有类 $cl(\mathcal{E},e)\in H^1_c(X\setminus S,\Pi^{\vee})$。记 $\ell_{S,\mathcal{E},\mathfrak{F}}$ 为一个带映射
\begin{equation}\tag{150}
(\ ,\ )_{\ell}:\phi_{S,\mathfrak{F}}\times\mathcal{E}_S\to\ell_{S,\mathcal{E},\mathfrak{F}}
\end{equation}
的 $\mu_\infty$-torsor，使得对 $\varphi\in\phi_{S,\mathfrak{F}}$、$e=(e_s)\in\mathcal{E}_S$、$h\in H^1(X\setminus S,\Pi(1))$、$\chi=(\chi_s)\in(\Pi^{\vee})^S$，有
\begin{equation}\tag{151}
\begin{aligned}
(\varphi+h,e)_{\ell}&=(h,cl(\mathcal{E},e))_P(\varphi,e)_{\ell},\\
(\varphi,\chi e)_{\ell}&=(\operatorname{Res}_S\varphi,\chi)(\varphi,e)_{\ell}.
\end{aligned}
\end{equation}
这里 $(\ ,\ )_P:H^1(X\setminus S,\Pi(1))\times H^1_c(X\setminus S,\Pi^{\vee})\to\mu_\infty$ 是 Poincaré 配对，而 $(\operatorname{Res}_S\varphi,\chi):=\prod_{s\in S}(\operatorname{Res}_s\varphi,\chi_s)\in\mu_\infty$。这样的 $(\ell_{S,\mathcal{E},\mathfrak{F}},(\ ,\ )_{\ell})$ 存在且唯一。若 $S'\supset S$，则有显然的映射 $\phi_{S,\mathfrak{F}}\hookrightarrow\phi_{S',\mathfrak{F}}$、$\mathcal{E}_{S'}\to\mathcal{E}_S$，且存在唯一的 $\mu_\infty$-torsor 等同 $\ell_{S,\mathcal{E},\mathfrak{F}}=\ell_{S',\mathcal{E},\mathfrak{F}}$，使这些映射关于 $(\ ,\ )_{\ell}$ 互为伴随。于是这个 $\mu_\infty$-torsor 不依赖 $S$，我们简记之为 $\ell_{\mathcal{E},\mathfrak{F}}$。

$\mathfrak{F}$ 变动时，$\ell_{\mathcal{E},\mathfrak{F}}$ 成为 $\operatorname{Bun}_G$ 上 $\mu_\infty$-torsor $\ell_{\mathcal{E}}$ 的纤维。函子
\begin{equation}\tag{152}
\ell=\ell^{\tilde G}:\Pi^{\vee}\operatorname{tors}(X)\to\mu_\infty\operatorname{tors}(\operatorname{Bun}_G),
\end{equation}
$\mathcal{E}\mapsto\ell_{\mathcal{E}}$，带有显然的 Picard 函子结构。相应的自同构群同态 $\Pi^{\vee}\to\Gamma(\operatorname{Bun}_G,\mu_\infty)$ 是 $\chi\mapsto(c,\chi)$。

\emph{注.} 事实上 $\ell$ 是函子 $\Pi^{\vee}\operatorname{tors}(X)\to\mu_m\operatorname{tors}(\operatorname{Bun}_G)$，其中 $m$ 是 $\Pi$ 的阶。这可由构造看出，也可由 (152) 是 Picard 函子推出。

\SourcePage{126}

\bdp{4.1.5} 对交换群 $A$，记 $A\operatorname{gerbes}(X)$ 为与 $X$ 上 $A$-gerbe 的 2-范畴相联系的范畴（故 $A\operatorname{gerbes}(X)$ 是群胚，其对象是 $X$ 上的 $A$-gerbe，其态射是 1-态射模去 2-同构）。4.1.2–4.1.4 实际上构造了双 Picard 函子
\begin{equation}\tag{153}
\Pi^{\vee}\operatorname{tors}(X)\times\Pi(1)\operatorname{gerbes}(X)\to\mu_\infty\operatorname{tors},
\end{equation}
其中 $\mu_\infty\operatorname{tors}$ 记一点上 $\mu_\infty$-torsor 的范畴。本小节（读者可以略过）给出该构造的「科学解释」。

在 [Del73] 的 §1.4.11 中，Deligne 给满足 $K^i=0$（$i\ne 0,-1$）的交换群复形 $K^{\cdot}$ 配置一个 Picard 范畴。记这个 Picard 范畴为 $P(K^{\cdot})$。其对象是 $K^0$ 的元素；从 $x\in K^0$ 到 $y\in K^0$ 的态射是满足 $df=y-x$ 的元素 $f\in K^{-1}$。

4.1.4 中我们隐含地把 $\Pi^{\vee}\operatorname{tors}(X)$ 解释为 $P(K^{\cdot}_S)$，其中 $K^0_S=H^1_c(X\setminus S,\Pi^{\vee})$ 是 $X$ 上在 $S$ 上平凡化的 $\Pi^{\vee}$-torsor 的同构类集合，$K^{-1}_S=H^0(S,\Pi^{\vee})$。4.1.3 中我们隐含地把 $\Pi(1)\operatorname{gerbes}(X)$ 解释为 $P(L^{\cdot}_S)$，其中 $L^0_S=H^2_S(X,\Pi(1))=\Pi^S$，$L^{-1}_S=H^1(X\setminus S,\Pi(1))$（$L^0_S$ 参数化 $X$ 上在 $X\setminus S$ 上固定了一个对象的 $\Pi(1)$-gerbe）。4.1.4 中给出的双 Picard 函子 (153) 的构造，只用到典范配对 $K^{\cdot}_S\times L^{\cdot}_S\to\mu_\infty[1]$。

对 $S'\supset S$ 有典范拟同构 $K^{\cdot}_{S'}\to K^{\cdot}_S$ 与 $L^{\cdot}_S\to L^{\cdot}_{S'}$。相应的范畴等价 $P(K^{\cdot}_{S'})\to P(K^{\cdot}_S)$ 与 $P(L^{\cdot}_S)\to P(L^{\cdot}_{S'})$，与我们关于 $P(K^{\cdot}_S)$、$P(K^{\cdot}_{S'})$ 同 $\Pi^{\vee}\operatorname{tors}(X)$ 的等同相容，也与关于 $P(L^{\cdot}_S)$、$P(L^{\cdot}_{S'})$ 同 $\Pi(1)\operatorname{gerbes}(X)$ 的等同相容。关于配对 $K^{\cdot}_S\times L^{\cdot}_S\to\mu_\infty[1]$ 与 $K^{\cdot}_{S'}\times L^{\cdot}_{S'}\to\mu_\infty[1]$，态射 $L^{\cdot}_S\to L^{\cdot}_{S'}$ 是 $K^{\cdot}_{S'}\to K^{\cdot}_S$ 的伴随。因此 (153) 不依赖 $S$。

\emph{注.}
% 「Remarks」在原书排在 p.126 末行，(i)(ii) 两项在 p.127；故页锚点放在标题之后。
\SourcePage{127}
\begin{enumerate}[label=(\roman*),leftmargin=2.2em]
\item 比起 $K^{\cdot}_S$ 与 $L^{\cdot}_S$，用它们在导出范畴中的像更自然，即 $(\tau_{\le 1}R\Gamma(X,\Pi^{\vee}))[1]$ 与 $(\tau_{\ge 1}R\Gamma(X,\Pi(1)))[2]$。但通常的导出范畴不够用：按 [Del73] 的 §§1.4.13–1.4.14，$K^{\cdot}$ 在导出范畴中的像只能把 $P(K^{\cdot})$ 确定到等价为止，而这个等价又只到相差一个同构为止，且该同构并不唯一。因此需要一个精细化了的导出范畴概念，它大概在文献中还找不到。
\item 由配对 $K^{\cdot}_S\times L^{\cdot}_S\to\mu_\infty[1]$ 非退化，容易推出 (153) 诱导 $\Pi^{\vee}\operatorname{tors}(X)$ 与 Picard 函子 $\Pi(1)\operatorname{gerbes}(X)\to\mu_\infty\operatorname{tors}$ 的范畴之间的等价（这是 [Del73] 中等价 (1.4.18.1) 的特例）。
\end{enumerate}

\bdp{4.1.6} 4.1.4 中 $\ell_{\mathcal{E}}$ 的定义可改写如下。设 $S\subset X$ 有限非空。对固定的 $e\in\mathcal{E}_S$，有类 $c=cl(\mathcal{E},e)\in H^1_c(X\setminus S,\Pi^{\vee})$，从而有态射 $\lambda_e:H^1(X\setminus S,\Pi(1))\to\mu_\infty$，由 $\lambda_e(h)=(h,c)_P$ 定义。记 $\ell_{\mathcal{E},e}$ 为 4.1.2 中 $H^1(X\setminus S,\Pi(1))$-torsor $\phi_S$ 沿 $\lambda_e$ 的前推。对不同 $e\in\mathcal{E}_S$ 的 torsor $\ell_{\mathcal{E},e}$，按下述方式等同。

设 $\tilde e=\chi e$，$\chi\in(\Pi^{\vee})^S$。则 $\lambda_{\tilde e}(h)/\lambda_e(h)=(\operatorname{Res}_S(h),\chi)$，其中 $\operatorname{Res}_S$ 是边缘态射 $H^1(X\setminus S,\Pi(1))\to H^2_S(X,\Pi(1))=\Pi^S$。于是 $\ell_{\mathcal{E},\tilde e}/\ell_{\mathcal{E},e}$ 是 $\Pi^S$-torsor $(\operatorname{Res}_S)_*\phi_S$ 沿 $\chi:\Pi^S\to\mu_\infty$ 的前推。4.1.3 中的映射 $\operatorname{Res}_{S,\mathfrak{F}}:\phi_{S,\mathfrak{F}}\to\Pi^S$ 诱导 $(\operatorname{Res}_S)_*\phi_S$ 的典范平凡化，从而诱导典范同构 $\ell_{\mathcal{E},e}\xrightarrow{\ \sim\ }\ell_{\mathcal{E},\tilde e}$。于是可把不同 $e\in\mathcal{E}_S$ 的 $\ell_{\mathcal{E},e}$ 等同起来，得到 $\operatorname{Bun}_G$ 上的 $\mu_\infty$-torsor，它不依赖 $e\in\mathcal{E}_S$。显然它也不依赖 $S$。这就是 $\ell_{\mathcal{E}}$。

\bdp{4.1.7} 设 $S\subset X$ 是非空有限集，$O_S:=\prod_{x\in S}O_x$，$K_S:=\prod_{x\in S}K_x$，其中 $O_x$ 是 $x$ 的完全化局部环，$K_x$ 是其分式域。记 $\underline{S}$ 为 $S$ 的形式邻域，$\operatorname{Bun}_{G,\underline{S}}$ 为 $X$ 上在 $\underline{S}$ 上平凡化的 $G$-丛的模概形（2.3.1 中引入的 $\operatorname{Bun}_{G,x}$ 对应 $S=\{x\}$）。在 $\operatorname{Bun}_{G,\underline{S}}$ 上定义 $G(K_S)$ 的作用，
% 原书 p.127 末句与 p.128 首句是同一句（"One defines an action ... extending the action of G(O_S) by interpreting ..."），
% 故此处不空行：页锚点落在同一段中间。
\SourcePage{128}
它延拓 $G(O_S)$ 的作用：把 $X$ 上的 $G$-丛解释为 $X\setminus S$ 上的 $G$-丛，连同其拉回到 $\operatorname{Spec} K_S$ 的一个平凡化（见 2.3.4 与 2.3.7）。

设 $\ell_{\mathcal{E}}$ 是 $\operatorname{Bun}_G$ 上对应于 $X$ 上 $\Pi^{\vee}$-torsor $\mathcal{E}$ 的 $\mu_\infty$-torsor（见 4.1.4、4.1.6）。记 $\ell^S_{\mathcal{E}}$ 为 $\ell_{\mathcal{E}}$ 在 $\operatorname{Bun}_{G,\underline{S}}$ 上的原像。$G(O_S)$ 在 $\operatorname{Bun}_{G,\underline{S}}$ 上的作用典范地提升为它在 $\ell^S_{\mathcal{E}}$ 上的作用。我们断言：$\mathcal{E}$ 在 $S$ 上的平凡化定义了 $G(K_S)$ 在 $\ell^S_{\mathcal{E}}$ 上的作用，它延拓上述 $G(O_S)$ 的作用，并与 $G(K_S)$ 在 $\operatorname{Bun}_{G,\underline{S}}$ 上的作用相容。事实上，一旦取定 $e\in\mathcal{E}_S$，就可把 $\ell^S_{\mathcal{E}}$ 与 $\ell^S_{\mathcal{E},e}=(\lambda_e)_*\tilde\phi_S$ 等同，其中 $\tilde\phi_S$ 是 $\phi_S$ 到 $\operatorname{Bun}_{G,\underline{S}}$ 的拉回，$\lambda_e$ 已在 4.1.6 中定义。$G(K_S)$ 作用在 $\tilde\phi_S$ 上，因为 $\phi_{S,\mathfrak{F}}$ 只依赖 $\mathfrak{F}$ 到 $X\setminus S$ 的限制。故 $G(K_S)$ 作用在 $\ell^S_{\mathcal{E},e}$ 上。

由 4.1.6 的同构 $\ell_{\mathcal{E},e}\xrightarrow{\ \sim\ }\ell_{\mathcal{E},\tilde e}$ 诱导的同构 $\ell^S_{\mathcal{E},e}\xrightarrow{\ \sim\ }\ell^S_{\mathcal{E},\tilde e}$ 并\emph{不}是 $G(K_S)$-等变的。事实上，若 $\tilde e=\chi e$，$\chi\in(\Pi^{\vee})^S$，则按 4.1.6，$\ell^S_{\mathcal{E},\tilde e}/\ell^S_{\mathcal{E},e}$ 是 $\Pi^S$-torsor $(\operatorname{Res})_*\tilde\phi_S$ 沿 $\chi:\Pi^S\to\mu_\infty$ 的前推。4.1.6 中的等同 $(\operatorname{Res})_*\tilde\phi_S=\operatorname{Bun}_{G,\underline{S}}\times\Pi^S$，在 $G(K_S)$ 经边缘态射 $\varphi:G(K_S)\to H^1(\operatorname{Spec} K_S,\Pi(1))=\Pi^S$ 作用在 $\Pi^S$ 上时，成为 $G(K_S)$-等变的（我们应当核对符号！！！）。因此平凡的 $\mu_\infty$-torsor $\ell^S_{\mathcal{E},\tilde e}/\ell^S_{\mathcal{E},e}$ 带上了 $G(K_S)$ 的非平凡作用：它按 $\chi\varphi:G(K_S)\to\mu_\infty$ 作用。

于是每个 $e\in\mathcal{E}_S$ 对应 $G(K_S)$ 在 $\tilde\phi_S$ 上的一个作用；若把 $e$ 换成 $\chi e$（$\chi\in(\Pi^{\vee})^S=\operatorname{Hom}(\Pi^S,\mu_\infty)$），则作用乘以 $\chi\varphi:G(K_S)\to\mu_\infty$。

\emph{注.} 顺便我们已证明：上边缘态射 $\varphi:G(K_S)\to H^1(\operatorname{Spec} K_S,\Pi(1))=\Pi^S$ 局部常值\footnote{另见 4.5.4。}（的确，$G(K_S)$ 作为群 ind-概形作用在 $(\operatorname{Res})_*\tilde\phi_S$ 上，故 $\varphi$ 是 ind-概形的态射，即 $\varphi$ 局部常值。证明可改述如下。不妨设 $S$ 只由一点 $x$ 组成。群 ind-概形 $G(K_x)$ 作用在 $\operatorname{Bun}_{G,x}$ 上（见 2.3.3--%
% 原书 p.128 末行与 p.129 首行是同一括注（"（see 2.3.3 –" / "2.3.4）"），故不空行。
\SourcePage{129}%
2.3.4），故它作用在 $\pi_0(\operatorname{Bun}_{G,x})=\pi_0(\operatorname{Bun}_G)$ 上。有「第一 Chern 类」映射 $c:\pi_0(\operatorname{Bun}_G)\to\Pi$。容易证明，对 $u\in\pi_0(\operatorname{Bun}_G)$、$g\in G(K_x)$，有 $c(gu)=\varphi(g)c(u)$，其中 $\varphi:G(K_x)\to H^1(K_x,\Pi(1))=\Pi$ 是上边缘态射。故 $\varphi$ 局部常值。

\bdp{4.1.8} 记 $\widetilde{G(K_S)_{\mathcal{E}}}$ 为由 $\mu_\infty$ 与元素 $\langle g,e\rangle$（$g\in G(K_S)$，$e\in\mathcal{E}_S$）生成的群，其定义关系为
\[
\begin{aligned}
\langle g_1g_2,e\rangle&=\langle g_1,e\rangle\langle g_2,e\rangle,\\
\langle g_1,\chi e\rangle&=\chi(\varphi(g))\cdot\langle g,e\rangle,\qquad \chi\in(\Pi^{\vee})^S=\operatorname{Hom}(\Pi^S,\mu_\infty),\\
\alpha\langle g,e\rangle&=\langle g,e\rangle\alpha,\qquad \alpha\in\mu_\infty.
\end{aligned}
\]
% [Tutor Note: 原书 p.129 第二条关系左端两处下标印作 g_1 与 g，不一致（左为 ⟨g_1,χe⟩，右为 χ(φ(g))·⟨g,e⟩）；疑右端下标应为 g_1，底本如此，保留原样。]

$\widetilde{G(K_S)_{\mathcal{E}}}$ 是 $G(K_S)$ 被 $\mu_\infty$ 的中心扩张。该扩张平凡：取定 $e\in\mathcal{E}_S$ 便给出分裂
\begin{equation}\tag{154}
\sigma_e:G(K_S)\to\widetilde{G(K_S)_{\mathcal{E}}},\qquad g\mapsto\langle g,e\rangle.
\end{equation}

由 4.1.7 可知，$\widetilde{G(K_S)_{\mathcal{E}}}$ 作用在 $\ell^S_{\mathcal{E}}$ 上：$\mu_\infty\subset\widetilde{G(K_S)_{\mathcal{E}}}$ 按显然方式作用，而对应于 $e\in\mathcal{E}_S$ 的 $G(K_S)$ 在 $\ell^S_{\mathcal{E}}$ 上的作用（见 4.1.7）来自分裂 (154)。

\bdp{4.1.9} 考虑 $\operatorname{Bun}_{G,\underline{S}}$ 中对应于 $X$ 上平凡 $G$-丛（连同其在 $\underline{S}$ 上的显然平凡化）的点。让 $G(K_S)$ 作用于此点，得到态射 $f:G(K_S)\to\operatorname{Bun}_{G,\underline{S}}$。设 $G$ 半单。则 $f$ 诱导同构
\begin{equation}\tag{155}
G(K_S)/G(A_S)\xrightarrow{\ \sim\ }\operatorname{Bun}_{G,\underline{S}},
\end{equation}
其中 $A_S:=H^0(X\setminus S,\mathcal{O}_X)$（见 [La-So] 的定理 1.3，以及该书 §3 中的证明）。关键在于把 $G(K_S)$ 与 $G(A_S)$ 看作群 ind-概形，并把 $G(K_S)/G(A_S)$ 看作 fppf 商；故 (155) 不止是 $\mathbb{C}$-点集合之间的双射。还有同构
\begin{equation}\tag{156}
G(O_S)\backslash G(K_S)/G(A_S)\xrightarrow{\ \sim\ }\operatorname{Bun}_G.
\end{equation}

\SourcePage{130}

容易看出，4.1.4 与 4.1.7 中定义的 $\mu_\infty$-torsor $\ell_{\mathcal{E}}$ 与 $\ell^S_{\mathcal{E}}$ 可描述为
\begin{equation}\tag{157}
\ell^S_{\mathcal{E}}=\widetilde{G(K_S)_{\mathcal{E}}}/G(A_S),
\end{equation}
\begin{equation}\tag{158}
\ell_{\mathcal{E}}=G(O_S)\backslash\widetilde{G(K_S)_{\mathcal{E}}}/G(A_S),
\end{equation}
\looseness=-1 其中 $\widetilde{G(K_S)_{\mathcal{E}}}$ 是 4.1.8 中的中心扩张。这里嵌入 $i:G(O)\to\widetilde{G(K_S)_{\mathcal{E}}}$ 与 $j:G(A_S)\to\widetilde{G(K_S)_{\mathcal{E}}}$ 由下式定义
\begin{equation}\tag{159}
i(g)=\langle g,e\rangle,\qquad e\in\mathcal{E}_S,
\end{equation}
\begin{equation}\tag{160}
j(g)=\langle g,e\rangle\cdot(\psi(g),cl(\mathcal{E},e))_P^{-1},\qquad e\in\mathcal{E}_S
\end{equation}
% [Tutor Note: 原书 p.130 此处印作 $i:G(O)\to\cdots$，与同一段 (158) 中的 $G(O_S)$ 不一致；疑为 $G(O_S)$ 之误，底本如此，保留原样。]
（我们应当核对符号！！！）其中 $\psi$ 是边缘态射 $G(A_S)\to H^1(X\setminus S,\Pi(1))$，$cl(\mathcal{E},e)\in H^1_c(X\setminus S,\Pi^{\vee})$ 是 $(\mathcal{E},e)$ 的类（(159) 与 (160) 的右端不依赖 $e$）。

\emph{注.} 态射 $\varphi:G(K_S)\to\Pi^S$ 与 $\psi:G(A_S)\to H^1(X\setminus S,\Pi(1))$ 诱导态射
\begin{equation}\tag{161}
\operatorname{Bun}_G=G(O_S)\backslash G(K_S)/G(A_S)\to\Pi^S/H^1(X\setminus S,\Pi(1)),
\end{equation}
其中 (161) 的右端理解为商叠。显然 $\ell_{\mathcal{E}}$ 是叠 $\Pi^S/H^1(X\setminus S,\Pi(1))$ 上某个 $\mu_\infty$-torsor 的拉回。

\bdp{4.1.10} 读者可以略过 4.1 的剩余部分。

设 $C$ 是群胚。记 $\underline{C}$ 为带 fppf 拓扑的 $\mathbb{C}$-概形范畴上对应的常群胚层。若 $C$ 中对象的自同构群有限，则 $\underline{C}$ 是代数叠。记号滥用之下，我们常把 $\underline{C}$ 写成 $C$（例如 $C$ 是集合时，$\underline{C}=C\times\operatorname{Spec}\mathbb{C}$ 常与 $C$ 等同）。

\emph{例.}
\begin{enumerate}[label=\arabic*),leftmargin=2.2em]
\item 若 $C$ 只有一个对象、$G$ 是其自同构群，则 $\underline{C}$ 是 $G$ 的分类叠。
\SourcePage{131}
\item 若 $C=P(K^{\cdot})$（见 4.1.5），则 $\underline{C}$ 是 $K^0$ 关于 $K^{-1}$ 作用的商叠。因此按 4.1.5，(161) 的右端是对应于群胚 $\Pi(1)\operatorname{gerbes}(X)$ 的叠。
\item 若 $C=A\operatorname{gerbes}(X)$，则 $\underline{C}$ 是与预层 $S\mapsto A\operatorname{gerbes}(X\times S)$ 相联系的群胚层。
\end{enumerate}

把群胚 $\Pi(1)\operatorname{gerbes}(X)$ 看作代数叠。4.1.2 中我们定义了典范态射
\begin{equation}\tag{162}
\tilde c:\operatorname{Bun}_G\to\Pi(1)\operatorname{gerbes}(X),
\end{equation}
它把 $G$-丛 $\mathfrak{F}$ 对应到 $\mathfrak{F}$ 的 $\tilde G$-提升构成的 $\Pi(1)$-gerbe（顺便说，对半单 $G$ 定义的态射 (161) 与 $\tilde c$ 一致）。$\tilde c$ 是 Chern 类映射 $c:\operatorname{Bun}_G\to H^2(X,\Pi(1))=\Pi$ 的精细化；更确切地说，$c$ 是 $\tilde c$ 与典范态射 $\Pi(1)\operatorname{gerbes}(X)\to H^2(X,\Pi(1))$ 的复合，其中 $H^2(X,\Pi(1))=\Pi$ 就是 $\Pi(1)\operatorname{gerbes}(X)$ 的同构类集合。

4.1.4 中构造的 $\operatorname{Bun}_G$ 上的 $\mu_\infty$-torsor，来自 $\Pi(1)\operatorname{gerbes}(X)$ 上的 $\mu_\infty$-torsor。下述命题表明：若 $\tilde G$ 是 $G$ 的泛覆叠，则 $\operatorname{Bun}_G$ 上任一局部系统都来自 $\Pi(1)\operatorname{gerbes}(X)$ 上唯一的局部系统。

\bdp{4.1.11} \emph{命题.} 设 $\tilde G$ 是 $G$ 的泛覆叠（故 $\Pi=\pi_1(G)$）。则态射 (162) 诱导 $\operatorname{Bun}_G$ 的基本群胚与 $\Pi(1)\operatorname{gerbes}(X)$ 之间的等价。

我们勾画一个超越的证明（既然是超越的，就不再区分 $\Pi$ 与 $\Pi(1)$）。记 $X^{\mathrm{top}}$ 为 $X$ 对应的 $C^\infty$ 流形；对 $X$ 上的 $G$-丛 $\mathfrak{F}$，记 $\mathfrak{F}^{\mathrm{top}}$ 为 $X^{\mathrm{top}}$ 上对应的 $G$-丛。考虑群胚 $\operatorname{Bun}_G^{\mathrm{top}}$，其对象是 $X^{\mathrm{top}}$ 上的 $G$-丛，其态射是 $G$-丛之间 $C^\infty$ 同构的同痕类。容易证明，自然函子 $\operatorname{Bun}_G^{\mathrm{top}}\to\Pi\operatorname{gerbes}(X^{\mathrm{top}})=\Pi\operatorname{gerbes}(X)$ 是等价。于是须证明：对 $X^{\mathrm{top}}$ 上的 $G$-丛 $\xi$，由 $X$ 上配备……的 $G$-丛 $\mathfrak{F}$ 构成的叠
% [Tutor Note: 原书 p.131 末行止于 "equipped"；该句的其余部分（"with an isotopy class F^top ≅ ξ, is non-empty, connected and simply connected"）在 p.132，见下一单元首行。译文在此以「……」标记断点。]
